\subsection*{日志重放面对两种应用状态}

让一个副本在应用 S1 后立即崩溃。它的内存余额、M 和编号计数器都消失，持久日志保留已提交的命令。副本从初始 10000 重建时，再应用 S1 得到 7000，这只是重建同一份本地状态。其他副本保存的 S1 也各自得到 7000。系统的逻辑扣款数量取决于状态机中的业务流水，三个副本保存的是同一笔扣款的三份镜像。

持久应用状态则有另一条故障历史。副本已经把余额 7000 和 D1 写入本地数据库，却在保存 \code{lastApplied=41} 以前崩溃。重启从旧边界 40 继续重放。如果 M 与 D1 已经共同保存，再应用 S1 可以根据 M 恢复原结果，并把应用边界推进。如果 M 也漏掉，重放就可能再次减款。

一种清楚的实现规则，是把应用后的业务数据、绑定和该次应用边界放在同一本地事务中。另一种是以内存状态为主，从完整快照和日志重建。两种实现都要给出崩溃后的唯一恢复起点。仅把 \code{lastApplied} 单独写到一个文件里，再假定它恰好与其他写入同时落盘，会留下混合状态。

\CElevenFigure{ch11-original12}{从空内存重建与在持久状态上继续重放，使用不同的恢复起点}{fig:ch11-replay-or-continue}

考虑另一方向：先持久保存边界 41，再保存余额和 M。崩溃后边界已经为 41，余额却还在 10000，恢复者会跳过 S1，造成遗漏。应用边界承担到这里为止的转换都已体现在这份状态中的含义。写大和写小都要与实际内容协调。

\CElevenFigure{ch11-applied-crashes}{先保存数据和先保存边界会留下两种相反错误，共同提交给出唯一恢复起点}{fig:ch11-applied-crashes}

前一节对业务身份去重，可以容忍同 S 位于多个日志位置。应用位置的协调还覆盖其他命令，包括调整配置、推进业务阶段和关闭会话。应用引擎把这些转换也纳入状态与位置的对应，每次恢复都从同一完整边界继续。

\subsection*{为什么快照需要整个应用状态}

日志不断增长，重启时从位置 1 重放的工作也随之增长。快照把应用到位置 k 以后的状态直接保存下来，后续恢复先安装它，再应用 k 以后的已提交后缀。日志截断指删除已经由可恢复快照覆盖的前缀，释放其存储与重放成本~\cite{raft}。

沿一条共同初态看这个问题：账户 1 原有 10000，E1 有效。S1/K1 成功 3000，S2/K2 又成功 3000。系统当前余额 4000，有 P1/O1/D1 与 P2/O2/D2，两份绑定和原结果。S1 的结果保存当次 7000，S2 保存当次 4000。设这些业务命令及后面的同参恢复，都已应用到日志位置 50。

快照若只复制当前余额 4000、两份订单与流水，重启后的入口就可能缺少 $M(S1)$。调用方重发 K1 时，受理入口看到有效 E1 和空绑定，按首次购买再次扣 3000，形成 P3/O3/D3，余额成为 1000。已有 D1 中或许保存着 S1，修复者可以设计由业务证据重建绑定的规则。这个错误入口直接开放重试，使保存在流水中的信息没有进入受理判断。

完整快照保存 W、M、$\Gamma$ 及其同一应用边界。选用协议序号适配的扩展状态另外保存 U，两个方向与 M 使用同一边界。M 保留 S1 的原输入、真实 P1 和原 $R(S1)$，恢复后同参命令直接得到原 P1/O1/D1 和 7000，当前余额继续为 4000。快照可以换一种等价编码保存这些信息，恢复所得判断与结果必须符合原规则。

\CElevenFigure{ch11-original13}{两种快照共用 4000 的业务初态，恢复入口遗漏绑定后旧 K1 会再次减款}{fig:ch11-complete-broken-snapshot}

$\Gamma$ 也有同样的作用。若 E1 已关闭，快照却把它恢复为有效，迟到旧请求便重新取得准入资格。若快照遗漏 $\Gamma$，而启动程序遇到未知 E1 就自动注册，同样绕过持久关闭。对已经受理的外部付款，原 C/T 及其未完成责任又要保留到服务端恢复终结。快照内容根据状态机承诺的恢复行为确定，当前余额只是其中一个分量。

\subsection*{用位置和任期给快照定位}

快照中的业务内容要能与日志接上。Raft 保存两个边界字段：\code{lastIncluded\allowbreak Index} 表示快照覆盖的最后日志位置 k，\code{lastIncluded\allowbreak Term} 表示该项原任期。它们为 k+1 的前驱检查提供依据。支持成员变更的系统还保存截至该边界生效的配置及相关位置，恢复者据此判断投票成员~\cite{raft}。

因此一份快照可以分为应用内容与复制元数据。应用内容是同一 $\Sigma_k$ 中的 W/M/$\Gamma$，复制元数据描述 k、该项任期以及所需配置。这里的 k 取自已提交、已应用的边界。领导者的日志尾部如果到了 55，提交边界为 52，应用只到 50，就只能把实际状态 50 与对应位置配在一起。

\CElevenFigure{ch11-original14}{快照内容和复制元数据共同定位到 k，恢复随后连接已提交日志后缀}{fig:ch11-snapshot-metadata}

把余额读于 50，M 读于 49，再贴上 k=50，是一种撕裂快照：各部分来自不同应用时点。仅在读取每张表时使用线程安全容器，保护的是局部访问，组合仍可能不一致。状态机可以短暂阻止应用，复制一份完整内存状态与边界，再恢复执行。写盘使用这份固定副本。也可以用写时复制或同一数据库快照，保持同一逻辑版本。写时复制让后续写入产生新版本，正在导出的旧版本继续指向原内容~\cite{raft,ongaro}。

例如导出者先取得 M 为空，应用线程随后完成 S1，导出者再读余额 7000。导出文件看起来包含正确余额，身份却来自更早时刻。完整状态复制要在应用串行边界完成，序列化到文件可以慢一些。先捕获一份逻辑状态，再把它写完，可以让应用继续推进而导出内容保持原版本。

\CElevenFigure{ch11-torn-snapshot}{逐项加锁也可能混合两个时点，先捕获完整状态再异步写盘可保住同一边界}{fig:ch11-torn-snapshot}

\subsection*{先完成可恢复副本，再删除被覆盖日志}

快照创建有自己的失败位置。导出临时文件时进程退出，可能只留下半个文件。把这个文件视为完整快照并删除前缀，恢复就缺少两份依据。持久快照路径先写完整内容、核对格式与完整性、完成所需刷新，再发布可用快照标识。只有成功发布并且元数据对应以后，才允许删除它覆盖的日志~\cite{paxos-made-live}。

日志已删、进程再崩溃，恢复从已发布快照开始。快照尚未发布就崩溃，恢复继续使用旧快照与尚未删除的日志。若应用状态存储和复制库分别掌握文件，就通过明确的完成通知和边界句柄连接：应用交回这份状态已保存至 k 的依据，复制层随后按 k 截断。《Paxos Made Live》为相同问题设计了快照句柄，包含日志实例位置与该点配置~\cite{paxos-made-live}。

\CElevenFigure{ch11-snapshot-lifecycle}{快照完成与日志截断之间的发布边界，决定每个崩溃点可用哪份恢复依据}{fig:ch11-snapshot-lifecycle}

落后副本可能需要一条已经被领导者删除的旧日志项。此时领导者发送完整快照，副本安装后再追赶后缀。安装期间对外读取和写入要看到受控状态，避免新 W 配上旧 M。收到的状态与本地日志在边界位置及任期匹配时，兼容后缀可以保留。其他情况由复制协议按安装规则重建对应日志。应用恢复过程沿库给出的边界整体替换旧状态~\cite{raft}。

第十章的查询投影快照使用对象业务版本 h/k，决定迟到投影能否向前安装。本节 Raft 快照使用复制日志位置与任期，覆盖整个应用状态。它们都要求状态与覆盖边界对应，比较对象和安装资格分别由各自协议定义。复制组安装一份备份时，依次核对配置、日志边界与前驱，再连接允许保留的后缀。

\subsection*{把领导者切换与快照恢复接在一起}

现在沿同一条封闭本地历史看最终行为。初始余额 10000，S1、S2 分别为 E1 下 K1、K2 的 3000 分购买。S1 应用并形成 P1/O1/D1、原 7000。成功响应在发送前丢失，原领导者退出。新领导者取得已提交前缀，重发的 S1 又进入日志，应用时读取 M，得到原结果。

S2 随后首次应用，形成 P2/O2/D2，当前余额 4000。S1 再到，服务继续返回原 7000。接着保存完整快照，确认它覆盖全部绑定、会话和业务状态以后截断前缀。各副本停止并恢复，先从快照得到当前 4000，再沿日志后缀追赶。S1 的下一次同参重试继续定位 P1。

\CElevenFigure{ch11-combined-history}{同一历史依次经过响应丢失、领导者切换、另一项购买与快照恢复}{fig:ch11-combined-history}

这个历史有三种不同计数。S1 可以多次被接收，多次进入已提交日志，多次调用应用函数。它在每个恢复所得逻辑状态中只关联一条 D1。进程重启可能再次调用应用函数以重建本地内存，或者在新位置处理同参恢复。实际购买效果由 M 与 W 中的关系计数。

将快照改为遗漏 M 的错误版本，其余执行保持相同，重启后便会出现第三笔扣款。复制层接受了一份按规定位置生成的应用快照，应用却少交了会改变后来请求解释的查找信息。恢复入口据此把原 S1 重新送入首次购买分支，错误沿完整复制流程传到各副本。

\subsection*{预测一项购买重试会得到什么}

设另一条历史从余额 10000 开始，E1 全程有效。S1/K1 成功 3000，得到 P1/O1/D1、原余额 7000。S2/K2 随后成功 3000，得到 P2/O2/D2、原余额 4000。当前余额 4000。完整状态已经复制并持久保留，领导者切换到拥有全部提交前缀的副本。

一位开发者提出缩小快照：保留全部购买、订单、流水和会话，另存一组 \code{seen\_keys}，表示哪些完整 S 已成功。恢复入口遇到已经出现的 S，就按 S 查出真实 P/O/D，再用当前余额构造成功响应。这里假定原输入也可以从成交记录完整取得，同参和异参比较各自正确。完整 $R(S)$ 没有单独保存。

请先预测三个结果：同参重试 S1 会得到哪张订单、哪个余额？这次重试是否增加扣款？如果紧接着重试 S2，两个结果又能否同时符合原接口？然后判断，把快照中的当前余额改成 S1 原来的 7000，能否修复这项设计。

原入口会让 S1 返回 O1 和当前 4000。订单身份正确，扣款数量保持两笔，返回的成交余额却已经改变。S2 恰好也返回 4000，看起来正确，是因为它在这条历史中最后执行。把当前余额改成 7000，又会使账户查询与两笔总计 6000 的支出不一致，并把 S2 的原结果改错。一个当前余额描述当前状态，两项成交则各有自己的历史余额。

\CElevenFigure{ch11-original-result}{当前余额与逐笔原结果分别保存，才能同时恢复 S1 的 7000 和 S2 的 4000}{fig:ch11-original-result}

可以保存完整 $R(S)$，也可以从不可变成交证据确定性重建原结果。本例若 D1/D2 分别保存当次余额、原报价、币种、金额、编号关系和固定截止等完整信息，恢复程序就有可能重建两项结果。程序在开放写入和重试以前完成重建，并验证 S、输入、P/O/D 与会话资格的关系，再把查询入口接到这份已恢复状态。如果流水已经足以承载全部必要信息，M 可以是一份索引。索引丢失时，恢复程序先完成明确的重建步骤，再开放依赖它的入口。

这也解释了前面的坏快照反例。它保留的流水可以含有旧身份和结果，故障来自恢复入口没有恢复或重建所需绑定，随后把旧 S1 当作首次购买。完整快照让入口在恢复后立即拥有原查找依据。

\subsection*{外部效果在同一恢复图中处于哪里}

把视线移回外部动作 \code{purchase.create.external.v2}。若复制的是购买服务应用状态，快照就要保留 A 固定的 S/P/C/T、平台范围和 \code{pay-C1}，R 已建立的 O1、C 中两类原结果、EV1，以及发布与核查的后续责任。P 继续只保存自己的业务进度。支付平台账本中的 X1 处于独立提交域。

假如三个副本在重放一条执行付款命令时直接调用外部平台，每个副本都可能重复发送。副本停止恢复后再重放，又会增加调用。正确组织是应用命令先保存相当于 A 的工作与固定身份，状态机外部发送者读取该事实，再按原 C1 与平台交互。可信结果回到另一个确定性命令，形成相当于 R 的结论。E 已经发生、R 尚未应用时，系统继续表达未决并查询原支付~\cite{rifl}。

最终拒绝的证据条件也随原协议保留：固定身份既往没有成功，所有已受理、在途和未来路径均已排除成功。一次查询没查到，或只观察到未来资格关闭时，原 C1 保持未决，恢复者继续核查其既往执行。已经复制的 \code{unconfirmed} 准确表达本地缺少终态证据，这种状态与平台已有 X1 可以同时存在。

\CElevenFigure{ch11-original15}{复制延续 A 与 R 的本地事实，平台 E 与状态机外发送者保持原支付身份}{fig:ch11-external-effects}

同理，EV1 重放可以使发布次数增加，履约库沿 I/B 恢复真实 F1。平台通知 PEV1 先被核验为 C1/X1 的证据，R 再建立本地 EV1。副本切换以后各段从自己的原身份继续：复制组用日志位置安排应用，平台用原付款身份判断重复，消费者用 I/B 恢复原接收与履约效果。
\FloatBarrier
